% Options for packages loaded elsewhere
\PassOptionsToPackage{unicode}{hyperref}
\PassOptionsToPackage{hyphens}{url}
\PassOptionsToPackage{space}{xeCJK}
\documentclass[
]{article}
\usepackage{xcolor}
\usepackage{amsmath,amssymb}
\setcounter{secnumdepth}{-\maxdimen} % remove section numbering
\usepackage{iftex}
\ifPDFTeX
  \usepackage[T1]{fontenc}
  \usepackage[utf8]{inputenc}
  \usepackage{textcomp} % provide euro and other symbols
\else % if luatex or xetex
  \usepackage{unicode-math} % this also loads fontspec
  \defaultfontfeatures{Scale=MatchLowercase}
  \defaultfontfeatures[\rmfamily]{Ligatures=TeX,Scale=1}
\fi
\usepackage{lmodern}
\ifPDFTeX\else
  % xetex/luatex font selection
  \ifXeTeX
    \usepackage{xeCJK}
    \setCJKmainfont[]{Microsoft YaHei}
  \fi
  \ifLuaTeX
    \usepackage[]{luatexja-fontspec}
    \setmainjfont[]{Microsoft YaHei}
  \fi
\fi
% Use upquote if available, for straight quotes in verbatim environments
\IfFileExists{upquote.sty}{\usepackage{upquote}}{}
\IfFileExists{microtype.sty}{% use microtype if available
  \usepackage[]{microtype}
  \UseMicrotypeSet[protrusion]{basicmath} % disable protrusion for tt fonts
}{}
\makeatletter
\@ifundefined{KOMAClassName}{% if non-KOMA class
  \IfFileExists{parskip.sty}{%
    \usepackage{parskip}
  }{% else
    \setlength{\parindent}{0pt}
    \setlength{\parskip}{6pt plus 2pt minus 1pt}}
}{% if KOMA class
  \KOMAoptions{parskip=half}}
\makeatother
\usepackage{color}
\usepackage{fancyvrb}
\newcommand{\VerbBar}{|}
\newcommand{\VERB}{\Verb[commandchars=\\\{\}]}
\DefineVerbatimEnvironment{Highlighting}{Verbatim}{commandchars=\\\{\}}
% Add ',fontsize=\small' for more characters per line
\newenvironment{Shaded}{}{}
\newcommand{\AlertTok}[1]{\textcolor[rgb]{1.00,0.00,0.00}{\textbf{#1}}}
\newcommand{\AnnotationTok}[1]{\textcolor[rgb]{0.38,0.63,0.69}{\textbf{\textit{#1}}}}
\newcommand{\AttributeTok}[1]{\textcolor[rgb]{0.49,0.56,0.16}{#1}}
\newcommand{\BaseNTok}[1]{\textcolor[rgb]{0.25,0.63,0.44}{#1}}
\newcommand{\BuiltInTok}[1]{\textcolor[rgb]{0.00,0.50,0.00}{#1}}
\newcommand{\CharTok}[1]{\textcolor[rgb]{0.25,0.44,0.63}{#1}}
\newcommand{\CommentTok}[1]{\textcolor[rgb]{0.38,0.63,0.69}{\textit{#1}}}
\newcommand{\CommentVarTok}[1]{\textcolor[rgb]{0.38,0.63,0.69}{\textbf{\textit{#1}}}}
\newcommand{\ConstantTok}[1]{\textcolor[rgb]{0.53,0.00,0.00}{#1}}
\newcommand{\ControlFlowTok}[1]{\textcolor[rgb]{0.00,0.44,0.13}{\textbf{#1}}}
\newcommand{\DataTypeTok}[1]{\textcolor[rgb]{0.56,0.13,0.00}{#1}}
\newcommand{\DecValTok}[1]{\textcolor[rgb]{0.25,0.63,0.44}{#1}}
\newcommand{\DocumentationTok}[1]{\textcolor[rgb]{0.73,0.13,0.13}{\textit{#1}}}
\newcommand{\ErrorTok}[1]{\textcolor[rgb]{1.00,0.00,0.00}{\textbf{#1}}}
\newcommand{\ExtensionTok}[1]{#1}
\newcommand{\FloatTok}[1]{\textcolor[rgb]{0.25,0.63,0.44}{#1}}
\newcommand{\FunctionTok}[1]{\textcolor[rgb]{0.02,0.16,0.49}{#1}}
\newcommand{\ImportTok}[1]{\textcolor[rgb]{0.00,0.50,0.00}{\textbf{#1}}}
\newcommand{\InformationTok}[1]{\textcolor[rgb]{0.38,0.63,0.69}{\textbf{\textit{#1}}}}
\newcommand{\KeywordTok}[1]{\textcolor[rgb]{0.00,0.44,0.13}{\textbf{#1}}}
\newcommand{\NormalTok}[1]{#1}
\newcommand{\OperatorTok}[1]{\textcolor[rgb]{0.40,0.40,0.40}{#1}}
\newcommand{\OtherTok}[1]{\textcolor[rgb]{0.00,0.44,0.13}{#1}}
\newcommand{\PreprocessorTok}[1]{\textcolor[rgb]{0.74,0.48,0.00}{#1}}
\newcommand{\RegionMarkerTok}[1]{#1}
\newcommand{\SpecialCharTok}[1]{\textcolor[rgb]{0.25,0.44,0.63}{#1}}
\newcommand{\SpecialStringTok}[1]{\textcolor[rgb]{0.73,0.40,0.53}{#1}}
\newcommand{\StringTok}[1]{\textcolor[rgb]{0.25,0.44,0.63}{#1}}
\newcommand{\VariableTok}[1]{\textcolor[rgb]{0.10,0.09,0.49}{#1}}
\newcommand{\VerbatimStringTok}[1]{\textcolor[rgb]{0.25,0.44,0.63}{#1}}
\newcommand{\WarningTok}[1]{\textcolor[rgb]{0.38,0.63,0.69}{\textbf{\textit{#1}}}}
\usepackage{longtable,booktabs,array}
\usepackage{calc} % for calculating minipage widths
% Correct order of tables after \paragraph or \subparagraph
\usepackage{etoolbox}
\makeatletter
\patchcmd\longtable{\par}{\if@noskipsec\mbox{}\fi\par}{}{}
\makeatother
% Allow footnotes in longtable head/foot
\IfFileExists{footnotehyper.sty}{\usepackage{footnotehyper}}{\usepackage{footnote}}
\makesavenoteenv{longtable}
\setlength{\emergencystretch}{3em} % prevent overfull lines
\providecommand{\tightlist}{%
  \setlength{\itemsep}{0pt}\setlength{\parskip}{0pt}}
\usepackage{bookmark}
\IfFileExists{xurl.sty}{\usepackage{xurl}}{} % add URL line breaks if available
\urlstyle{same}
\hypersetup{
  hidelinks,
  pdfcreator={LaTeX via pandoc}}

\author{}
\date{}

\begin{document}

\section{面向噪声标签数据的细粒度图像识别鲁棒微调 ---
技术报告}\label{ux9762ux5411ux566aux58f0ux6807ux7b7eux6570ux636eux7684ux7ec6ux7c92ux5ea6ux56feux50cfux8bc6ux522bux9c81ux68d2ux5faeux8c03-ux6280ux672fux62a5ux544a}

\begin{quote}
赛题：全球校园人工智能算法精英大赛 · 算法挑战赛（初赛阶段）
骨干网络：CLIP ViT-B/32（OpenAI 官方预训练权重，timm
\texttt{vit\_base\_patch32\_clip\_224.openai}） 日期：2026 年 6 月 12 日
\end{quote}

\begin{center}\rule{0.5\linewidth}{0.5pt}\end{center}

\subsection{一、算法概述}\label{ux4e00ux7b97ux6cd5ux6982ux8ff0}

本方案针对含噪标签条件下的细粒度图像识别任务，提出一套''\textbf{特征空间自动去噪
+ 软标签自训练 + 参数高效鲁棒微调}''的三段式微调框架。

首先，利用冻结的 CLIP ViT-B/32 提取全部训练图像的视觉特征，基于
\textbf{kNN
邻域标签一致性}对每个样本的标注可信度进行无监督评分：若一个样本的近邻多数与其标签一致，则该标注大概率正确。据此按类别自适应地保留各类前
75\%
高一致性样本，并对一致性极高的样本设置保留保底，避免误删少数难样本。其次，在筛选后的子集上训练余弦分类器作为教师，对被过滤样本中模型预测置信度高者赋予\textbf{软伪标签}回收利用，实现''过滤不浪费''。最后，在视觉骨干的全部注意力层注入
\textbf{LoRA 低秩适配器}（仅 1.27M 可训练参数，约占骨干参数的
1.4\%），骨干原始权重完全冻结，分类头由教师模型热启动，以软标签加权损失进行端到端微调。低秩约束从结构上限制了表征漂移，兼顾噪声鲁棒性与预训练知识保留，天然抑制灾难性遗忘。

由于训练标签含噪、无干净验证集，本方案还设计了\textbf{按邻域一致性分带的本地验证协议}：以''标签基本正确但样本困难''的中间一致性带（0.3--0.6）准确率作为核心选型指标，有效规避了噪声验证集''奖励拟合噪声''的评估偏差。在本地留出集上，最终模型相对无约束基线全集准确率提升约
9--15 个百分点，整套流程（含噪声筛选）全自动、固定随机种子、一键复现。

\subsection{二、实现方案}\label{ux4e8cux5b9eux73b0ux65b9ux6848}

\subsubsection{2.1 解题思路}\label{ux89e3ux9898ux601dux8def}

赛题数据具有三个核心难点：标签噪声重、类间差异细微（细粒度）、无干净验证集。我们的总体思路是：

\begin{enumerate}
\def\labelenumi{\arabic{enumi}.}
\tightlist
\item
  \textbf{不依赖人工清洗}------噪声甄别必须自动化、可复现（满足赛题第五条算法设计要求）；
\item
  \textbf{过滤要保守、回收要充分}------细粒度场景下''低置信样本''中混有大量''难而正确''的样本，激进过滤会损伤表征（实验证实
  keep 50\%/60\% 全面劣于 75\%）；
\item
  \textbf{微调要受约束}------以 LoRA
  低秩适配替代全参数微调，从参数空间结构上限制对预训练知识的破坏；
\item
  \textbf{评估要去偏}------构建分带验证协议解决''用噪声标签评估模型''的悖论。
\end{enumerate}

\subsubsection{2.2 数据处理}\label{ux6570ux636eux5904ux7406}

\begin{itemize}
\tightlist
\item
  \textbf{数据概况}：500 个细粒度类别，103218 张训练图（类内 201--213
  张，分布均衡），24967
  张测试图；类别目录为纯数字编号，无语义类名，故未使用文本模态先验。
\item
  \textbf{损坏容错}：部分图片文件截断，统一开启
  \texttt{PIL\ ImageFile.LOAD\_TRUNCATED\_IMAGES}，全量图片可正常读取。
\item
  \textbf{特征缓存}：冻结骨干对全量图片单次推理（含水平翻转 TTA
  取平均、L2 归一化），缓存 768
  维特征，使后续所有去噪与头部实验在秒级完成，支撑了高密度消融。
\item
  \textbf{噪声扫描}：kNN（k=16，余弦相似度加权）邻域标签一致性统计显示，42\%
  训练样本一致性 \textless{} 0.3，21\% \textless{}
  0.1，证实噪声水平较高。
\end{itemize}

\subsubsection{2.3
算法设计与开发}\label{ux7b97ux6cd5ux8bbeux8ba1ux4e0eux5f00ux53d1}

\textbf{第一阶段：kNN 邻域一致性选样。}
对每个训练样本，在全体训练特征中检索 16
个最近邻（排除自身），以相似度加权计算近邻标签与自身标签的一致率。按类别取一致性前
75\% 样本，同时无条件保留一致性 ≥ 0.7
的样本。按类自适应选样保证了筛选不会系统性偏向某些类别。最终保留
77494/103218（75.08\%）。

\textbf{第二阶段：教师训练与软伪标签回收。}
在保留子集上训练余弦分类器（特征与权重均 L2
归一化、可学习温度，对噪声更稳健）。训练完成后对被过滤样本推理，预测置信度
≥ 0.7 者以预测类别构造平滑软标签、按置信度加权重新纳入训练（标签平滑
0.1）。

\textbf{第三阶段：LoRA 鲁棒微调。} 在 ViT 全部 12 层注意力的 qkv
与输出投影线性层注入 LoRA（秩 r=16，α=32，dropout
0.05），骨干原权重全部冻结；分类头由第二阶段教师热启动。训练目标为软标签交叉熵的样本加权平均。数据增强采用
RandomResizedCrop(0.65--1.0) + 水平翻转 + 轻度
ColorJitter（噪声数据下避免重增强）。

\subsubsection{2.4
模型训练与优化}\label{ux6a21ux578bux8badux7ec3ux4e0eux4f18ux5316}

\begin{itemize}
\tightlist
\item
  优化器 AdamW，LoRA 参数 lr 2e-4（无权重衰减）、分类头 lr 1e-3（wd
  1e-2），OneCycle 调度（10\% warmup），混合精度训练，batch size 192；
\item
  10 个 epoch，单卡 RTX 5060 Laptop（8GB）每 epoch 约 5 分钟，总训练约
  50 分钟，工程成本低；
\item
  推理采用单模型 + 水平翻转
  TTA（单一推理流程，符合''禁止多模型集成''约束）。
\end{itemize}

\subsubsection{2.5 测试与验证}\label{ux6d4bux8bd5ux4e0eux9a8cux8bc1}

\textbf{分带验证协议}（核心方法论贡献）：按类分层留出 10\%（10285
张）作验证集，并按验证样本自身的 kNN 一致性分为三带：

\begin{itemize}
\tightlist
\item
  低带（\textless0.3，4308
  张）：标签多为错标，\textbf{此带准确率高反而意味着拟合噪声}；
\item
  \textbf{中带（0.3--0.6，2673
  张）：标签基本正确但样本难，区分度最高，作为模型选型主指标}；
\item
  高带（≥0.6，3304 张）：接近饱和（98\%+），仅作健康检查。
\end{itemize}

\textbf{消融实验结果}（同一划分、同一随机种子）：

\begin{longtable}[]{@{}llll@{}}
\toprule\noalign{}
方案 & 噪声全集 & 中带 & 高带 \\
\midrule\noalign{}
\endhead
\bottomrule\noalign{}
\endlastfoot
置信度两阶段选样（复现的对照基线） & 0.5840 & --- & --- \\
纯 CE + 标签平滑（无去噪） & 0.6409 & 0.8058 & 0.9897 \\
GMM 损失建模加权 & 0.5934 & --- & --- \\
kNN 过滤（keep 50\%） & 0.5928 & 0.7415 & 0.9840 \\
kNN 过滤（keep 75\%）+ 软标签自训练 & 0.6450 & 0.8182 & 0.9906 \\
\textbf{+ LoRA 微调（最终方案）} & \textbf{0.7353} & \textbf{0.8887} &
\textbf{0.9918} \\
\end{longtable}

最终方案相对无去噪基线：噪声全集 +9.4pt、中带
+8.3pt；相对复现的置信度选样基线全集
+15.1pt。提交文件经自动化脚本校验（文件名与测试集严格一致、四位类别编号、zip
规范）。

\subsection{三、算法创新点}\label{ux4e09ux7b97ux6cd5ux521bux65b0ux70b9}

\begin{enumerate}
\def\labelenumi{\arabic{enumi}.}
\tightlist
\item
  \textbf{kNN 邻域一致性去噪}：完全在冻结 CLIP
  特征空间进行，无需训练即可对 10 万级样本在 2
  秒内完成噪声评分；按类自适应阈值 +
  高一致保底的''保守过滤''设计，避免了细粒度场景下难样本被误删------消融显示这是过滤策略成败的关键。
\item
  \textbf{``过滤不浪费''的软伪标签回收}：被过滤样本经教师模型高置信重标注后以软标签、置信度加权方式回收，在去噪与数据利用率之间取得平衡。
\item
  \textbf{分带验证协议}：针对''无干净验证集''这一赛题核心困境，提出按邻域一致性分带评估、以中带准确率选型的方法，识别并规避了噪声验证集奖励噪声拟合的系统性偏差（实验中多个''全集指标更高''的配置在低带暴露出噪声拟合迹象而被否决）。
\item
  \textbf{结构化抗遗忘}：LoRA 低秩约束 + 骨干冻结 +
  教师头热启动三重机制保留预训练先验，全部可训练参数仅 1.27M，在 8GB
  消费级显卡上即可完成训练，方案对复赛/半决赛更大规模数据具备良好的可扩展性。
\end{enumerate}

\subsection{四、问题与思考}\label{ux56dbux95eeux9898ux4e0eux601dux8003}

\begin{longtable}[]{@{}
  >{\raggedright\arraybackslash}p{(\linewidth - 2\tabcolsep) * \real{0.5000}}
  >{\raggedright\arraybackslash}p{(\linewidth - 2\tabcolsep) * \real{0.5000}}@{}}
\toprule\noalign{}
\begin{minipage}[b]{\linewidth}\raggedright
遇到的问题
\end{minipage} & \begin{minipage}[b]{\linewidth}\raggedright
分析与解决
\end{minipage} \\
\midrule\noalign{}
\endhead
\bottomrule\noalign{}
\endlastfoot
噪声验证集上''不去噪的模型分数更高''，去噪方法看似全部失效 &
定位为评估偏差：噪声验证集奖励拟合噪声的模型。设计分带验证协议，以中带为主指标后各方法优劣立判 \\
激进过滤（保留 50--60\%）反而全面掉点 &
细粒度任务中低一致性样本混有大量难而正确的样本；改为保守过滤（75\%）+
高一致保底 + 软标签回收 \\
GMM 损失建模、多轮自训练等经典方法收益不及预期 & GMM
在细粒度+均衡分布下区分度不足；多轮自训练低带准确率持续抬升（拟合噪声迹象），均未采用------经验：方法选择须以去偏指标为准，不迷信文献默认配置 \\
Windows 下 DataLoader 验证阶段 worker 进程崩溃 & 系统内存不足（每 worker
为独立进程约 1GB）；验证 loader 降为 2 worker 并持久化复用 \\
EMA + mixup 组合掉点 & 训练总步数少（\textasciitilde180 步），EMA decay
0.999 远未收敛；短程头部训练不适用，弃用 \\
部分图片文件截断报错 & 按赛题提示开启 Pillow 截断容错，全量可读 \\
\end{longtable}

\textbf{进一步优化方向}：① 全量数据（含验证留出）重训最终模型；②
延长训练（epoch 10 仍为 best，未见过拟合）；③ 以微调后特征迭代刷新 kNN
噪声评分（特征-去噪交替优化）；④ 复赛长尾场景下引入 logit adjustment
与类频感知的选样阈值；⑤
取得线上分数后标定本地中带指标与线上准确率的相关性，降低提交消耗。

\subsection{五、过程进度}\label{ux4e94ux8fc7ux7a0bux8fdbux5ea6}

\begin{longtable}[]{@{}
  >{\raggedright\arraybackslash}p{(\linewidth - 6\tabcolsep) * \real{0.2500}}
  >{\raggedright\arraybackslash}p{(\linewidth - 6\tabcolsep) * \real{0.2500}}
  >{\raggedright\arraybackslash}p{(\linewidth - 6\tabcolsep) * \real{0.2500}}
  >{\raggedright\arraybackslash}p{(\linewidth - 6\tabcolsep) * \real{0.2500}}@{}}
\toprule\noalign{}
\begin{minipage}[b]{\linewidth}\raggedright
阶段
\end{minipage} & \begin{minipage}[b]{\linewidth}\raggedright
开始
\end{minipage} & \begin{minipage}[b]{\linewidth}\raggedright
结束
\end{minipage} & \begin{minipage}[b]{\linewidth}\raggedright
主要完成内容
\end{minipage} \\
\midrule\noalign{}
\endhead
\bottomrule\noalign{}
\endlastfoot
数据准备 & 2026-06-11 & 2026-06-11 &
数据下载解压、完整性扫描、类别分布统计、冻结特征全量提取与缓存 \\
基线构建 & 2026-06-11 & 2026-06-12 & 复现冻结特征 + 余弦头 +
置信度两阶段选样基线，首版提交文件生成 \\
算法设计与消融 & 2026-06-12 & 2026-06-12 & 建立分带验证协议；两轮共 17
组消融（kNN/GMM/软标签/mixup/EMA/keep-ratio/多轮自训练），锁定
kNN75+软标签配方 \\
模型训练与优化 & 2026-06-12 & 2026-06-12 & LoRA
微调管线开发、冒烟测试、10 epoch 完整训练（中带 0.8887） \\
测试与提交 & 2026-06-12 & 2026-06-12 & 测试集推理（翻转
TTA）、提交文件生成与自动化格式校验 \\
文档撰写 & 2026-06-12 & 2026-06-12 &
实验记录（EXPERIMENTS.md）与本技术报告 \\
\end{longtable}

\subsection{六、团队分工}\label{ux516dux56e2ux961fux5206ux5de5}

\begin{longtable}[]{@{}lll@{}}
\toprule\noalign{}
成员 & 角色 & 负责内容 \\
\midrule\noalign{}
\endhead
\bottomrule\noalign{}
\endlastfoot
（待填写） & 算法设计 & 鲁棒微调框架设计、去噪策略与验证协议设计 \\
（待填写） & 编程开发 & 训练/推理管线实现、实验执行、工程优化 \\
（待填写） & 数据分析 & 噪声特征分析、消融实验结果分析 \\
（待填写） & 文档撰写 & 技术报告、实验记录、提交材料整理 \\
\end{longtable}

\subsection{七、解题参考资源}\label{ux4e03ux89e3ux9898ux53c2ux8003ux8d44ux6e90}

\begin{enumerate}
\def\labelenumi{\arabic{enumi}.}
\tightlist
\item
  Radford A, et al.~Learning Transferable Visual Models From Natural
  Language Supervision. ICML 2021.（CLIP）
\item
  Hu E J, et al.~LoRA: Low-Rank Adaptation of Large Language Models.
  ICLR 2022.
\item
  Li J, et al.~DivideMix: Learning with Noisy Labels as Semi-supervised
  Learning. ICLR 2020.（损失建模选样思想）
\item
  Guo Y, Gu X. JoAPR: Cleaning the Lens of Prompt Learning for
  Vision-Language Models. CVPR 2024.
\item
  Zhang X, et al.~TrustCLIP: Learning from Noisy Labels via Semantic
  Label Verification and Trust-aligned Gradient Projection. ACM MM 2025.
\item
  Wortsman M, et al.~Robust Fine-tuning of Zero-shot Models (WiSE-FT).
  CVPR 2022.（鲁棒微调与抗遗忘思想）
\item
  工具：PyTorch、timm、Pillow、NumPy；预训练权重：OpenAI CLIP
  ViT-B/32（经 HuggingFace timm 发布的官方权重）。
\end{enumerate}

\begin{center}\rule{0.5\linewidth}{0.5pt}\end{center}

\subsubsection{附：复现命令}\label{ux9644ux590dux73b0ux547dux4ee4}

\begin{Shaded}
\begin{Highlighting}[]
\ExtensionTok{pip}\NormalTok{ install }\AttributeTok{{-}r}\NormalTok{ requirements.txt}
\ExtensionTok{python}\NormalTok{ main.py                                   }\CommentTok{\# 特征提取（生成 outputs/cache）}
\ExtensionTok{python}\NormalTok{ exp\_head.py                               }\CommentTok{\# 第一轮消融}
\ExtensionTok{python}\NormalTok{ exp\_round2.py                             }\CommentTok{\# 第二轮调优}
\ExtensionTok{python}\NormalTok{ finetune\_lora.py }\AttributeTok{{-}{-}epochs}\NormalTok{ 10              }\CommentTok{\# LoRA 微调（90/10 验证）}
\ExtensionTok{python}\NormalTok{ finetune\_lora.py }\AttributeTok{{-}{-}predict}                \CommentTok{\# 生成 pred\_results\_lora.csv/zip}
\ExtensionTok{python}\NormalTok{ check\_submission.py }\AttributeTok{{-}{-}csv}\NormalTok{ pred\_results\_lora.csv }\AttributeTok{{-}{-}zip}\NormalTok{ pred\_results\_lora.zip}
\end{Highlighting}
\end{Shaded}

全流程固定随机种子（seed=42），噪声筛选环节完全由代码自动执行，无人工清洗步骤。

\end{document}
